\documentclass[11pt,fleqn]{article}
% \documentclass[smallextended]{svjour3}

\textheight=23.3cm 
\textwidth=16.5cm 
\topmargin=-15mm 
\oddsidemargin=0cm 

\usepackage{enumerate}
\usepackage{amsmath}
\usepackage{amssymb}
\usepackage{bbm}
\usepackage{ulem}
\usepackage{graphicx}
%This package is dated and it is giving me problems
%\usepackage{subfig}
\usepackage{wrapfig}
\usepackage{pdflscape}
\usepackage[sidewaysfigure]{rotating}
\usepackage{natbib}
%\usepackage{amsthm}
%\usepackage{setspace}
\usepackage{multirow}
\newcommand{\tb}{\hspace*{3mm}}
\bibpunct{(}{)}{;}{a}{,}{,}
\parskip=3mm 
\parindent=5mm 
\baselineskip=17pt

\usepackage{amsthm}
\newtheorem{assumption}{Assumption}
\newtheorem{definition}{Definition}
\newtheorem{proposition}{Proposition}
\newtheorem{corollary}{Corollary}
\newtheorem{lemma}{Lemma}
\newtheorem{theorem}{Theorem}
\renewcommand{\qedsymbol}{\textit{Q.E.D.}}

\DeclareMathOperator*{\argmax}{arg\,max}
\DeclareMathOperator*{\argmin}{arg\,min}

\usepackage{algpseudocode}
\usepackage{algorithm}
\renewcommand\algorithmicthen{}
\renewcommand\algorithmicdo{}

\usepackage{hyperref}
\hypersetup{
  pdftitle={Semi-Endogenous Shocks},    % title
  pdfauthor={Robert Kirkby},     % author
  %hyperfootnotes=false,
  %hyperindex=true,
  %hyperfigures=true,  
  colorlinks=true,                % makes links a coloured word, rather than putting coloured boxes around them
  linkcolor=black,                % makes internal links black
  citecolor=black,                % makes citations black
  urlcolor=blue,                 % color of external links
  breaklinks=true                 % allow links to be brocken across multiple lines      
}


\begin{document}


%===============================
\newcommand{\rob}[1]{{\noindent\color{blue}\bf #1}}
\newcommand{\robc}[1]{{\noindent\color{blue}\bf R: #1}}
%===============================

\title{VFI Toolkit: Pseudocodes}
\author{Robert Kirkby\thanks{Kirkby: Victoria University of Wellington.
Please address all correspondence about this article to Robert Kirkby at $<$robertdkirkby@gmail.com$>$,
\href{http://www.robertdkirkby.com}{robertdkirkby.com}, \href{http://www.vfitoolkit.com}{vfitoolkit.com}.
}} % \\ {\small Victoria University of Wellington}}
% \author{Robert Kirkby \thanks{I thank}}
% \institute{Victoria University of Wellington, \email{robertdkirkby@gmail.com},
% \href{http://www.robertdkirkby.com}{robertdkirkby.com}.}
\date{\today}
\maketitle

% \noindent {\bf Keywords:} Value function iteration, Semi-endogenous shocks.

% \noindent {\bf JEL Classification:} E00; C68; C63; C62

\thispagestyle{empty}

%==============================================================================
\newpage
{
\small
\parskip=0cm
\parindent=0cm
\tableofcontents
}
%==============================================================================


\newpage
\baselineskip=17pt
\setcounter{page}{1}

\section{Introduction}
This document presents pseudocodes for various commands from VFI Toolkit. The actual implementations in VFI Toolkit often involve parallelized versions of the psuedocodes presented here, and also typically have options to use more-or-less parallelization (versus looping).

This document is a work-in-progress. Please request pseudocode for any specific command you want to see it for (either by email or on discourse.vfitoolkit.com).

Throughout I use the standard notation of VFI Toolkit: so $d$ is a decision variable, $a$ is endogenous state (and $a'$ the next-period endogenous state), $z$ is a markov exogenous state, $e$ is an i.i.d. exogenous state, $u$ is an i.i.d. shock that occurs between periods. The transition probability from current exogenous state $z$ to next-period exogenous state $z'$ is denoted by $\pi_z(z,z')$. $V$ is the value function. For finite horizon problems there are $N_j$ periods, indexed by $j$.

Throughout, $D(\cdot)$ denotes the set of feasible choices conditional on the current states.

In many psuedocodes $g(a,z)$ denotes a policy, with $g^d(a,z)$ being the policy for the decision variables and $g^{a'}(a,z)$ being the policy for the next-period endogenous state.

\section{Value Function Iteration: Solving for V and Policy}

\subsection{Infinite Horizon: decision variable, markov shock}
Starting from an initial guess for the value function, we iterate on the value function until convergence. To speed things up Howards improvement is used.
\begin{algorithmic}
 \State Declare initial value $V_0$.
 \State Declare iteration count $n=0$.
 \While {$||V_n-V_{n-1}||>$ 'tolerance constant'}
   \State Increment $n$. Let $V_{old}=V_{n-1}$.
   \For All values $z$
    \State Calculate $E[V_{old}(a',z')]$
    \For All values of $a$
     \State Calculate $V_n(a,z)=max_{d,a' \in D(a,z)} F(d,a',a,z)+\beta E[V_{old}(a',z')]$
     \State ... and keep the $argmax$ $g_n(a,z)$.
     \EndFor
   \EndFor
   \If UseHowards==1
    \For nHowards number of times
     \State Update $V_n(a,z)=F(g^{d}_n(a,z),g^{a'}_n(a,z),a,z)+\beta E[V_n(g^{a'}_n(a,z),z')]$
    \EndFor
   \EndIf
 \EndWhile
 \\ \Return $V_n$, $g_n$
\end{algorithmic}
Note that by making the outer-loop over $z$ we are able to reduce the number of times we compute the expectations (which depend on $z$ and not on $a$).

Howards improvement algorithm is to use the current policy to evaluate/update the value function a couple of times. This is faster as it skips the maximization step (which is computationally the most costly) and because the policy function typically is 'pointing' in the direction of the solution, so a few cheap computations to move further in this direction is worthwhile. $UseHowards$ is implemented as conditions under which Howards improvement algorithm should be used (which can be controlled using $vfoptions$. $nHowards$ is set to 80 by default (based on a trying out a bunch of different values, this seemed to be best at speeding up the 'average' problem). In practice, we avoid using Howards improvement algorithm for the first few iterations (as the initial guess is likely poor, and so Howards does not much help), and also stop using Howards improvement algorithm once we are within one order of magnitude of convergence of the value function.\footnote{This is because otherwise Howards leads to very small errors. You won't notice them normally, but when trying to compute a placebo transition path between stationary equilibrium they are problematic as the transition won't use Howards and so tries to go to a very slightly different place. If you are thinking 'but I saw the proof in Sargent \& Stachurski that Howards is fine and gives same answer' the trick is that the last line of their proof assumes that Howards evaluates the 'policy greedy' value function, but in practice you will never do this. (Placebo transition=a transition path in which nothing changes, this is actually quite a demanding test of accuracy; this is not a dig at S\&S, all textbooks do the same.)}

A couple of hints to how this could be done better (the kind of improvements that work for just about any algorithm, rather than things like using endogenous grid method instead of pure discretization): ideally it would use the improved versions of Howards developed in \cite{Rendahl2022}, \cite{BakotaKredler2022} and \cite{PhelanEslami2022} (I've not tried these). And it would use relative VFI, or even endogenous VFI (for these later two, they were tried but don't work in VFI Toolkit because of the pure discretization);  \cite{Bray2017}.

\subsubsection{Refinement} \label{sec:refine}
If you set \textit{vfoptions.solnmethod = ‘purediscretization\_{}refinement’;} then the algorithm will be the 'refinement'. Refinement is based on the observation that $d$ does not enter the expectations next period value function, only the return function. So we can 'pre-solve',  choosing $d$ to maximize the return function $F(d,aprime,a,z)$ to get $d^{*}(aprime,a,z)$, the decision that maximizes the return function given $(aprime,a,z)$, and associated maximum values $F^{*}(aprime,a,z)$. We then perform standard value function iteration, but just using $F^{*}(aprime,a,z)$, so we have essentially removed $d$ from the problem on which we have to iterate. Once we finish solving we can put the optimal policy for $aprime$ into $d^{*}(aprime,a,z)$ to get the optimal policy for $d$. This is faster because we just pre-solve for $d$ once, and thereby avoid having to solve for it as part of every iteration.

\begin{algorithmic}
   \For All values $z$  \Comment First, pre-solve for $d$
    \For All values $a$
     \State Calculate $F^{*}(a',a,z)=max_{d \in D(a,z)} F(d,a',a,z)$ and the argmax $d^{*}(a',a,z)$
    \EndFor
   \EndFor
 \State Declare initial value $V_0$. \Comment Now do standard VFI without $d$ variable
 \State Declare iteration count $n=0$.
 \While {$||V_n-V_{n-1}||>$ 'tolerance constant'}
   \State Increment $n$. Let $V_{old}=V_{n-1}$.
   \For All values $z$
    \State Calculate $E[V_{old}(a',z')]$
    \For All values of $a$
     \State Calculate $V_n(a,z)=max_{a' \in D(a,z)} F^{*}(a',a,z)+\beta E[V_{old}(a',z')]$
     \State ... and keep the $argmax$ $g_n(a,z)$.
     \EndFor
   \EndFor
   \If UseHowards==1
    \For nHowards number of times
     \State Update $V_n(a,z)=F^{*}(g^{a'}_n(a,z),a,z)+\beta E[V_n(g^{a'}_n(a,z),z')]$
    \EndFor
   \EndIf
 \EndWhile
 \State Evaluate $g_n^{d}(a,z)=d^{*}(g^{a'}(a,z),a,z)$  \Comment Recover the policy for $d$
 \\ \Return $V_n$, $g_n$
\end{algorithmic}

Note that the only difference between the refinement is internal. All inputs and outputs will be identical. Refinement also has the nice property that if we loop when calculating $d^{*}$ and then parallize the value function iteration itself, we reduce the memory requirements as well as reducing the runtime making it possible to solve more complex problems (as long as they have a $d$ variable). Obviously the refinement is not possible in models that do not have a $d$ variable to begin with.\footnote{Pure-discretization makes this refinement easy, but in principle it should be possible to do with algorithms that use functional forms to fit the value and policy functions.}



\subsection{Finite Horizon: decision variable, markov shock}
The basic algorithm is simple backward iteration. Let $\Theta$ be the full set of model parameters, and let $\theta_j$ be the values of those parameters that are relevant in period $j$.\footnote{So if a parameter is independent of age, then it appears unchanged in $\theta_j$. If a parameter depends on age, then $\theta_j$ contains only the value relevant at age $j$.}
\begin{algorithmic}
 \State Solve for final period, $N_j$: $V_{N_j}(a,z)=\max_{d,a' \in D(a,z)} F(d,a',a,z) $
 \State ... and keep the $\argmax$ $g_{N_j}$.
 \For $j$ count backward from $N_j-1$ to $1$
   \State Get $\theta_j$ from $\Theta$
   \For All values of $z$
    \State Calculate $E[V_{j+1}(a',z')|z]=\sum_{z'} V_{j+1}(a',z') \pi_z(z,z')$
    \For All values of $a$
     \State Calculate $V_j(a,z)=\max_{d,a' \in D(a,z)} \left\{ F(d,a',a,z)+\beta E[V_{j+1}(a',z')|z] \right\}$
     \State ... and keep the $\argmax$ $g_j(a,z)$.
     \EndFor
   \EndFor
  \EndFor
   \\ \Return $[V_1,V_2,...,V_{N_j}]$, $[g_1,g_2,...,g_{N_j}]$
\end{algorithmic}

If you were writing custom code for a specific problem, you would likely try to exploit features such as monotonicity. However, because VFI Toolkit is designed to solve generic problems, it uses this simple and robust approach.

\subsubsection{Finite Horizon: grid interpolation layer}
Grid interpolation layer adds a second layer of points for the next-period endogenous state. The first layer is the original grid on $a'$, while the second layer places $n2$ evenly spaced points between adjacent first-layer grid points. This follows closely the pseudocode in \textit{GridInterpLayer.pdf}, but here we describe the case with one endogenous state, one decision variable, and one markov shock, matching the standard finite-horizon problem above.

Let $a\_{}grid=\{a_1,\dots,a_n\}$ be the grid on assets. The fine grid used by the second layer is
\[
aprime\_{}grid=\texttt{interp1}(1\!:\!1\!:\!n,\;a\_{}grid,\;\texttt{linspace}(1,n,n+n2*(n-1))).
\]
This is the grid that contains the original $n$ first-layer points together with the $n2$ evenly spaced points between each adjacent pair of first-layer points.

\begin{algorithmic}
 \State For final period, $N_j$, set $E[V_{N_j+1}(a',z')|z]=0$
 \For $j$ count backward from $N_j$ to $1$
   \State Get $\theta_j$ from $\Theta$
   \For All values of $z$
    \If{$j<N_j$}
     \State Calculate $E[V_{j+1}(a',z')|z]=\sum_{z'} V_{j+1}(a',z') \pi_z(z,z')$
    \Else
     \State $E[V_{j+1}(a',z')|z]=0$
    \EndIf
    \State Interpolate $E[V_{j+1}(a',z')|z]$ from $a\_{}grid$ onto $aprime\_{}grid$ to get $\widehat{E[V_{j+1}]}(\tilde{a}',z)$
    \For All values of $a$
     \Statex \textit{First Layer}
     \For All values of $d$
      \State Solve for the optimal $a'$ on the first-layer grid:
      \Statex \hspace{\algorithmicindent} $g_j^1(d,a,z)=\argmax_{i=1,\dots,n} \left\{ F(d,a\_{}grid(i),a,z)+\beta E[V_{j+1}(a\_{}grid(i),z')|z] \right\}$
      \State Let $midpoint_j(d,a,z)=\max(\min(g_j^1(d,a,z),n-1),2)$
     \EndFor
     \Statex \textit{Second Layer}
     \For All values of $d$
      \State Create the second-layer grid around the first-layer midpoint:
      \Statex \hspace{\algorithmicindent} $aprimeindexes_j(d,a,z)=(midpoint_j+(midpoint_j-1)*n2)+(-n2-1:1:1+n2)$
      \Statex \hspace{\algorithmicindent} $aL2\_{}grid_j(d,a,z)=aprime\_{}grid(aprimeindexes_j(d,a,z))$
     \EndFor
     \State Solve for the optimal $(d,a')$ on the second layer:
     \Statex \hspace{\algorithmicindent} $(d_j^*(a,z),g_j^2(a,z))$
     \Statex \hspace{\algorithmicindent} $=\argmax_{d,\;i=1,\dots,3+2*n2} \left\{ F(d,aL2\_{}grid_j(d,a,z,i),a,z)+\beta \widehat{E[V_{j+1}]}(aL2\_{}grid_j(d,a,z,i),z) \right\}$
     \State Set $V_j(a,z)$ equal to this maximum.
    \EndFor
   \EndFor
  \EndFor
  \State Convert the first-layer midpoint and second-layer index into the stored policy.
  \State Let the lower-grid-point index be $g_{j,lower}(a,z)=midpoint_j(d_j^*(a,z),a,z)-\mathbbm{1}\{g_j^2(a,z)<n2+2\}$.
  \State Let the interpolation-layer index be
  \Statex \hspace{\algorithmicindent} $g_{j,L2}(a,z)=\mathbbm{1}\{g_j^2(a,z)<n2+2\}g_j^2(a,z)+\mathbbm{1}\{g_j^2(a,z)\geq n2+2\}(g_j^2(a,z)-n2-1)$.
  \State Store the policy as $g_j(a,z)=[d_j^*(a,z),g_{j,lower}(a,z),g_{j,L2}(a,z)]$.
   \\ \Return $[V_1,V_2,...,V_{N_j}]$, $[g_1,g_2,...,g_{N_j}]$
\end{algorithmic}

The first layer is essentially the usual discretized problem, except that with a decision variable it is solved conditional on each $d$. The second layer then creates a different local fine grid for each possible value of $d$, and the final maximization over the second layer chooses the optimal $(d,a')$. This is exactly how the Matlab code in \texttt{ValueFnIter\_{}FHorz\_{}GI\_{}raw.m} is organized.

The policy is stored using the lower grid point on the original $a\_{}grid$ together with the interpolation-layer index, rather than storing the value of $a'$ directly. This makes it easy to use the policy later for both the agent distribution and evaluation on the agent distribution. The upper grid point is just one plus the lower grid point, and because the second-layer points are evenly spaced, the interpolation weights can be recovered directly from the interpolation-layer index.

\subsubsection{Finite Horizon: semi-exogenous state}
We now add a semi-exogenous state. Semi-exogenous shocks are exogenous states for which the transition probabilities depend on a decision variable (see Life-Cycle Model 30 for an example). Thus, $\pi_{semiz}(semiz'|semiz,d)$ gives the probability of the next-period semi-exogenous state as a function of the current semi-exogenous state and the decision variable.

The main difference is that the expectation term is now more complex.

\begin{algorithmic}
 \State Solve for final period, $N_j$: $V_{N_j}(a,semiz,z)=\max_{d,a' \in D(a,semiz,z)} F(d,a',a,semiz,z)$
 \State ... and keep the $\argmax$ $g_{N_j}$.
  \For $j$ count backward from $N_j-1$ to $1$
   \State Get $\theta_j$ from $\Theta$
   \For All values of $d$
    \For All values of $(semiz,z)$
     \State Calculate $E[V_{j+1}(a',semiz',z')|d,semiz,z]$
     \Statex \hspace{\algorithmicindent} $=\sum_{semiz',z'} V_{j+1}(a',semiz',z') \pi_z(z,z') \pi_{semiz}(semiz'|semiz,d)$
     \For All values of $a$
      \State Calculate $V_j(a,semiz,z|d)$
      \Statex \hspace{\algorithmicindent} $=\max_{a' \in D(a,semiz,z)} \left\{ F(d,a',a,semiz,z)+\beta E[V_{j+1}(a',semiz',z')|d,semiz,z] \right\}$
      \State ... and keep the $\argmax$ $g_j(a,semiz,z|d)$.
      \EndFor
     \EndFor
    \EndFor
    \For All values of $(a,semiz,z)$
     \State $V_j(a,semiz,z)=\max_{d} V_j(a,semiz,z|d)$
     \State Let $d^*(a,semiz,z)=\argmax_d V_j(a,semiz,z|d)$.
     \State Then $g_j(a,semiz,z)=g_j(a,semiz,z|d^*(a,semiz,z))$.
    \EndFor
  \EndFor
   \\ \Return $[V_1,V_2,...,V_{N_j}]$, $[g_1,g_2,...,g_{N_j}]$
\end{algorithmic}

The term $E[V_{j+1}(a',semiz',z')|d,semiz,z]$ is doing a lot of work here. It adds an extra dimension to the expectation term, because the expectation must now be computed separately for each value of $d$ and for each current pair $(semiz,z)$. For this reason, it is natural to solve the problem conditional on $d$ first, and then maximize over $d$ afterward. Here, values of $d$ that are not in $D(a,semiz,z)$ are implicitly assigned a return of $-\infty$, so they will never be chosen.


\subsubsection{Finite Horizon: experienceasset}
'experienceasset' is an endogenous state where aprime cannot be chosen directly, but is instead a function of a decision variable and this period endogenous state; $aprime(d,a)$. The basic idea of how VFI Toolkit handles these is to evaluate $aprime(d_i,a_i)$ as a value on the grids of d and a, and then linearly interpolate this $aprime(d_i,a_i)$ onto the grid for a by allocating it to the grid-points above and below that value with linear weights.

So in a finite-horizon model with no other endogenous states, and no decision variables except those influencing the next period experienceasset, the algorithm is
\begin{algorithmic}
 \State Solve for final period, $N_j$: $V_{N_j}(a,z)=max_{d\in D(a,z)} F(d,a,z)$
 \State ... and keep the argmax $g_{N_j}$.
 \For $j$ count backward from $N_j-1$ to $1$
   \State Get $\theta_j$ from $\Theta$
   \For All values $z$
    \State Calculate $E[V_{j+1}(a',z')]$
    \State Evaluate $aprime(d,a)$ on the grids for $d$ and $a$
    \State Interpolate $aprime(d,a)$ onto the grid for $a$; get $aprime_i(d,a)$ and $aProbs(d,a)$, the index for the lower grid point and the probability of the lower grid point.
    \State Switch $E[V_{j+1}(a',z')]$ to $E[V_{j+1}(d,a,z')]$ using $aprime_i(d,a)$ and $aProbs(d,a)$ to replace $a'$ with $d$ and $a$.
    \For All values of $(a,z)$
     \State Calculate $V_j(a,z)=max_{d \in D(a,z)} F(d,a,z)+\beta E[V_{j+1}(d,a,z')]$
     \State ... and keep the $argmax$ $g_j(a,z)$.
     \EndFor
   \EndFor
  \EndFor
   \\ \Return $[V_1,V_2,...,V_{N_j}]$, $[g_1,g_2,...,g_{N_j}]$
\end{algorithmic}
notice that $a'$ no longer appears in the return function. The main trick is to replace $E[V_{j+1}(a',z')]$ to $E[V_{j+1}(d,a,z')]$, by interpolating $aprime(a,z)$ onto the grid on $a$.

The way interpolation is done is that for a value of $aprime$ that we want to interpolate onto the grid $[a_1,a_2,\dots,a_n]$ we, (i) find the $a_i$ that is the largest grid point that has a value less than (or equal to) $aprime$, call this the lower grid point, (ii) we allocate $aprime$ to both the lower grid point $a_i$ and the upper grid point $a_{i+1}$ using weights, (iii) we use linear weights for $a_i$ and $a_{i+1}$, so if $aprime$ is halfway between the two we give a weight of 0.5 to each, this is implemented by giving the lower grid point, $a_i$ a weight of $1-(aprime-a_i)/(a_{i+1}-a_i)$ (the weight for the upper grid point is just 1-this). This interpolation can be seen in \href{https://github.com/vfitoolkit/VFIToolkit-matlab/blob/master/aprimeFnMatrix/CreateExperienceAssetFnMatrix_Case1.m}{CreateExperienceAssetFnMatrix\_{}Case1.m}, the first part of the code creates $aprime(a,z)$ and then the last 30-40ish lines are doing the interpolation onto the grid, note that we only need to keep the lower grid point index and the probability assigned to the lower grid point (as the upper grid point is just adding one to this index, and the probability is one minus this).

\subsubsection{How interpolation is done}\label{sec:interp}
VFI Toolkit uses linear interpolation as part of alternative endogenous states like experienceasset, experienceassetu, riskyasset, and more. There are two halves to how the toolkit uses interpolation. One half is to put points onto a grid (as a lower-grid-point index and the interpolation weight on that lower grid point, which together implicitly determine the upper grid point and its weight). The second half is to use this to 'switch' grids.

For example in the previous pseudocode we saw the first half as: \\
\textit{Interpolate $aprime(d,a)$ onto the grid for $a$; get $aprime_i(d,a)$ and $aProbs(d,a)$, the index for the lower grid point and the interpolation weight on the lower grid point.}
\\ The way that the toolkit interpolates a value onto the grid is to use linear interpolation. So say we have that the value we want to interpolate onto the grid is $x$, and the grid is $\{x_1,x_2,\dots,x_n\}$. If $x\leq x_1$ we assign all the weight to $x_1$. If $x\geq x_n$ we assign all the weight to $x_n$. Otherwise, we find the two grid points either side of $x$, call them $x_i$ and $x_{i+1}$, so that $x_i\leq x < x_{i+1}$. Then we call $x_i$ the lower grid point, and assign it the interpolation weight $1-\frac{(x-x_i)}{(x_{i+1}-x_i)}$ (which is between 0 and 1, and is 1 if $x=x_i$). The interpolation weight on the upper grid point is then $\frac{(x-x_i)}{(x_{i+1}-x_i)}$, which is also between 0 and 1 and approaches 1 as $x$ approaches $x_{i+1}$. We only need to store the lower-grid-point index and the interpolation weight on the lower grid point, because the upper grid point is then just the next point on the grid and its weight is one minus the lower-grid-point weight.

While the second half in the previous pseudocode we saw as: \\
\textit{Switch $E[V_{j+1}(a',z')]$ to $E[V_{j+1}(d,a,z')]$ using $aprime_i(d,a)$ and $aProbs(d,a)$ to replace $a'$ with $d$ and $a$.}
\\ The way the toolkit does this is that for every $(d,a)$ we have a value $aprime(d,a)$ which we associated in the first half with a lower-grid-point index and interpolation weight, $aprime_i(d,a)$ and $aProbs(d,a)$. Looking at $E[V_{j+1}(a',z')]$ and ignoring the $E[\cdot]$ for the moment, we get
\begin{eqnarray*}
V_{j+1}(a',z') &=& aProbs(d,a)\,V_{j+1}(a_{aprime_i(d,a)},z') \\
&& + (1-aProbs(d,a))\,V_{j+1}(a_{aprime_i(d,a)+1},z') \\
&\equiv& V_{j+1}(d,a,z').
\end{eqnarray*}
At the lower and upper boundaries this collapses to putting all weight on a single grid point. We just need to do this before we evaluate the expectations term, and we will thus get $E[V_{j+1}(d,a,z')]$.

\subsubsection{Finite Horizon: riskyasset}
'riskyasset' is an endogenous state where aprime cannot be chosen directly, but is instead a function of a decision variable and a between-period i.i.d. shock; $aprime(d,u)$. The basic idea of how VFI Toolkit handles these is to evaluate $aprime(d_i,u_i)$ as a value on the grids of d and u, and then linearly interpolate this $aprime(d_i,u_i)$ onto the grid for $a$ by allocating it to the grid-points above and below that value with linear weights.

So in a finite-horizon model with no other endogenous states, and no decision variables except those influencing the next period riskyasset, the algorithm is
\begin{algorithmic}
 \State Solve for final period, $N_j$: $V_{N_j}(a,z)=max_{d\in D(a,z)} F(d,a,z)$
 \State ... and keep the argmax $g_{N_j}$.
 \For $j$ count backward from $N_j-1$ to $1$
   \State Get $\theta_j$ from $\Theta$
   \For All values $z$
    \State Calculate $E[V_{j+1}(a',z')]$
    \State Evaluate $aprime(d,u)$ on the grids for $d$ and $u$
    \State Interpolate $aprime(d,u)$ onto the grid for $a$; get $aprime_i(d,u)$ and $aProbs(d,u)$, the index for the lower grid point and the probability of the lower grid point.
    \State Switch $E[V_{j+1}(a',z')]$ to $E[V_{j+1}(d,u,z')]$ using $aprime_i(d,u)$ and $aProbs(d,u)$ to replace $a'$ with $d$ and $u$.
    \State Take expectations over $u$ to convert $E[V_{j+1}(d,u,z')]$ to $E[V_{j+1}(d,z')]$
    \For All values of $(a,z)$
     \State Calculate $V_j(a,z)=max_{d \in D(a,z)} F(d,a,z)+\beta E[V_{j+1}(d,z')]$
     \State ... and keep the $argmax$ $g_j(a,z)$.
     \EndFor
   \EndFor
  \EndFor
   \\ \Return $[V_1,V_2,...,V_{N_j}]$, $[g_1,g_2,...,g_{N_j}]$
\end{algorithmic}
notice that $a'$ no longer appears in the return function. The main trick is to replace $E[V_{j+1}(a',z')]$ to $E[V_{j+1}(d,z')]$, by interpolating $aprime(d,u)$ onto the grid on $a$, and then taking expectations over $u$.

Note: An important improvement to runtimes in models with riskyasset is the use of 'refinement'. Essentially in many models some decision variables will appear in the ReturnFn but not the aprimeFn, and vice-versa. We can 'refine' out these decision variables before we combine the return function, $F(d,a,z)$ and the expectations term, $E[V_{j+1}]$. This massively shrinks the size of the combined problem and leads to both faster runtimes and lower memory use. It only makes sense to do on a gpu. Refine is explained for infinite horizon problems in Section \ref{sec:refine}. The linear interpolation is explained in Section \ref{sec:interp}.

\subsubsection{Finite Horizon: standard \& riskyasset}
When using a standard endogenous state and a riskyasset together, the risky asset is always set up as the second endogenous state so we will call them $a1$ and $a2$, respectively. This approach is just to combine the usual approaches to each endogenous state, so the standard endogenous state will have $a1prime$ chosen on the grid while the riskyasset will be evaluated on the grids of $d$ and $u$ to get $a2prime(d_i,u_i)$ and then this is linearly interpolated onto the $a2$ grid by allocating it to the grid-points above and below that value with linear weights.

So in a finite-horizon model with no decision variables except those influencing the next period riskyasset, the algorithm is
\begin{algorithmic}
 \State Solve for final period, $N_j$: $V_{N_j}(a1,a2,z)=max_{a1,d\in D(a1,a2,z)} F(d,a1prime,a1,a2,z)$
 \State ... and keep the argmax $g_{N_j}$.
 \For $j$ count backward from $N_j-1$ to $1$
   \State Get $\theta_j$ from $\Theta$
   \For All values $z$
    \State Calculate $E[V_{j+1}(a1',a2',z')]$
    \State Evaluate $a2prime(d,u)$ on the grids for $d$ and $u$
    \State Interpolate $a2prime(d,u)$ onto the grid for $a2$; get $a2prime_i(d,u)$ and $a2Probs(d,u)$, the index for the lower grid point and the probability of the lower grid point.
    \State Switch $E[V_{j+1}(a1',a2',z')]$ to $E[V_{j+1}(d,a1',u,z')]$ using $a2prime_i(d,u)$ and $a2Probs(d,u)$ to replace $a2'$ with $d$ and $u$.
    \State Take expectations over $u$ to convert $E[V_{j+1}(d,a1',u,z')]$ to $E[V_{j+1}(d,a1',z')]$
    \For All values of $(a1,a2,z)$
     \State Calculate $V_j(a1,a2,z)=max_{d,a1' \in D(a1',a2',z)} F(d,a1',a1,a2,z)+\beta E[V_{j+1}(d, a1',z')]$
     \State ... and keep the $argmax$ $g_j(a1,a2,z)$.
     \EndFor
   \EndFor
  \EndFor
   \\ \Return $[V_1,V_2,...,V_{N_j}]$, $[g_1,g_2,...,g_{N_j}]$
\end{algorithmic}
notice that $a2'$ does not appear in the return function. The main trick is to replace $E[V_{j+1}(a1',a2',z')]$ with $E[V_{j+1}(d,a1',z')]$, by interpolating $a2prime(d,u)$ onto the grid on $a2$, and then taking expectations over $u$.

\subsubsection{Finite Horizon: Epstein-Zin preferences}
There are various different ways to set up Epstein-Zin preferences, such as in consumption-units, in utility-units with positive utility, in utility-units with negative utility, and 'double Epstein-Zin' which treats the additional risk aversion differently for survival probabilities to other probabilities. VFI Toolkit handles all of these in a single code by taking advantage of the fact that we can set up a generic Epstein-Zin with eight 'Epstein-Zin coefficients' (called $ezc1$, $ezc2$,$\dots$, $ezc8$ in the codes), and then all the different Epstein-Zin preferences simply correspond to different values of these eight coefficients.\footnote{This approach was developed from scratch for VFI Toolkit.}

The generic Epstein-Zin value fn problem for a finite-horizon value fn problem with markov shocks is,
\begin{eqnarray*}
 V_j(a,z)=\max_{d,a'} ezc1*\left( (ezc1*F(d,a',a,z)^{ezc2}) + ezc3*\beta*(s_j E[ezc4*V_{j+1}(a',z')^{ezc5}|z]^{ezc8})^{ezc6} \right)^{ezc7}
\end{eqnarray*}

The algorithm for solving this is,
\begin{algorithmic}
 \State Solve for final period, $N_j$: $V_{N_j}(a,z)=max_{d,a' \in D(a,z)} (ezc1*F(d,a',a,z)^{ezc2})^{ezc7} $
 \State ... and keep the argmax $g_{N_j}$.
 \For $j$ count backward from $N_j-1$ to $1$
   \State Get $\theta_j$ from $\Theta$
   \For All values $z$
    \State Calculate $E[ezc4*V_{j+1}(a',z')^{ezc5}|z]^{ezc8}$
    \For All values of $(a,z)$
     \State Calculate $V_j(a,z)=max_{d,a' \in D(a,z)} ezc1*\left( (ezc1*F(d,a',a,z)^{ezc2})+ezc3*\beta (s_j E[ezc4*V_{j+1}(a',z')^{ezc5}|z]^{ezc8})^{ezc6} \right)^{ezc7}$
     \State ... and keep the $argmax$ $g_j(a,z)$.
     \EndFor
   \EndFor
  \EndFor
   \\ \Return $[V_1,V_2,...,V_{N_j}]$, $[g_1,g_2,...,g_{N_j}]$
\end{algorithmic}
Implementing the code for this requires a bit of care because there are issues around things like 0s and negative powers where you have to be careful to treat it correctly.

We now go through different Epstein-Zin setups and explain how the eight Epstein-Zin coefficients are set in each case. We use $ceis$ and $crisk$ to refer to the elasticity of intertemporal substitition and risk aversion parameters, respectively.
\\ \noindent \textit {Traditional EZ in consumption units}
\begin{eqnarray*}
ezc1=1; \\
ezc2=1-1./ceis; \\
ezc3=1; \\
ezc4=1; \\
ezc5=1-crisk; \\
ezc6=(1-1./ceis)./(1-crisk); \\
ezc7=1./(1-1./ceis); \\
ezc8=1;
\end{eqnarray*}
\\ \noindent \textit{EZ in utility-units (in utils): positive utility}
\begin{eqnarray*}
 ezc1=1; \\
 ezc2=1; \\
 ezc3=1; \\
 ezc4=1; \\
 ezc5=1-crisk; \\
 ezc6=1./(1-crisk); \\
 ezc7=1; \\
 ezc8=1;
\end{eqnarray*}
\\ \noindent \textit{EZ in utility-units (in utils): negative utility}
\begin{eqnarray*}
 ezc1=1; \\
 ezc2=1; \\
 ezc3=-1; \\
 ezc4=-1; \\
 ezc5=1+crisk; \\
 ezc6=1./(1+crisk); \\
 ezc7=1; \\
 ezc8=1;
\end{eqnarray*}
Note: If the utility is negative use 1+crisk instead of 1-crisk. This way the interpretation of crisk is identical in both cases. And then we need to multiply it by -1 in two places, which is done by setting $ezc3$ and $ezc4$ to both be negative.
\\ \noindent \textit{Double Epstein-Zin: inner EZ is about risk, out EZ is about mortality risk}
\begin{eqnarray*}
\text{EZ in consumption units:} \\
    ezc6=(1-1./ceis)./(1-mrisk); \\
    ezc8=(1-mrisk)./(1-crisk); \\
\text{EZ in utils with positive utility: } \\
    ezc6=1./(1-mrisk); \\
    ezc8=(1-mrisk)./(1-crisk); \\
\text{EZ in utils with negative utility: } \\
    ezc6=1./(1+mrisk); \\
    ezc8=(1+mrisk)./(1+crisk);
\end{eqnarray*}
Note: this must be used together with one of the other settings, hence just showing which parameters need to change. $mrisk$ is the risk aversion parameter for the mortality risk.
\\ \noindent \textit{Scale returns by $1-\beta$ or $1-s_j\beta$:}
Some people like to scale the return function by $1-\beta$ or $1-s_j\beta$ (done with vfoptions.EZoneminusbeta). This just requires setting $ezc1=1-\beta$ or $ezc1=1-s_j\beta$.

There is one aspect of Epstein-Zin preferences that cannot be accomodated by this approach, namely when the elasticity of intertemporal substitution is equal to one. As can clearly be seen in the definitions of the Epstein-Zin coefficients, this would lead to some of them being undefined due to division by zero, or blowing off to infinity. It is possible to solve Epstein-Zin preferences with the elasticity of intertemporal substitution is equal to one, but it cannot be fit into our generic Epstein-Zin problem. For how the problem has to be set up, see Sargent \& Lungquist 4ed (2018) on pages 566-568, specifically looking at eqn (14.7.3) and compare the two cases. In practice, you could use the toolkit to solve with an elasticity of intertemporal substitition equal to 0.99 and again with 1.01, and as long as they give much the same solution, just interpret this as being a good enough approximation to the solution with 1.\footnote{Thanks to Louis Daumas for pointing me to where to find this.}



\section{Agent Distribution}

\subsection{Stationary Agent Distribution in Infinite Horizon}
There are three was to compute the stationary agent distribution in infinite horizon models. The stationary agent definition is defined as $\mu(a,z)$ satisfying $\mu(a',z')=P(a',z'|a,z) \mu(a,z)$, where $P$ is the markov transition kernel that combines the policy function $g(a'|a,z)$ with the exogenous shock transition kernel $\pi_z(z'|z)$ ($P=g \circ \pi_z$).

The three ways to compute the stationary agent distribution are (i) as the eigenvector of $\mu = P \mu$, (ii) by simulation, and (iii) by iteration.

Which of the three to use? The iteration method is the default as it has a better combination of speed and accuracy than the simulation method. However the iteration method does require more memory and so for large models it is not an option; if you are getting out of memory errors when solving for the agent distribution then turn off iteration by setting \textit{simoptions.iterate=0} and the toolkit will instead use the simulation. The eigenvecter method is unstable so is never used.

\subsubsection{Stationary Agent Distribution as Eigenvector}
In principle this is simple, we know the stationary distribution is defined as satisfying $\mu = P \mu$, and once the problem is discretized $P$ is a matrix. So $\mu$ is just an eigenvector of the matrix $P$ and there are plenty of routines existing for finding eigenvectors in all standard software and they are very fast. In practice however the method is unstable as minor numerical rounding errors in computing the eigenvector mean it often contains negative elements. There are ways to 'bend' the eigenvector so that it is a probability distribution (all elements are positive, sum to one) as given, e.g., 'Step 5' on page 178 of \cite*{BadshahBeaumontSrivastava2013} describes how to find the normalized eigenvector of the unit eigenvalue of $P$ and use the Perron-Frobenius theorem to 'bend' it.

By default the toolkit will avoid the eigenvector approach due to the instability (if you want to use the eigenvector method, set \textit{simoptions.eigenvector=1}. Neither of the ways to 'bend' it have been implemented (if you want to contribute one, please do :)

\subsubsection{Simulate Agent Distribution}
We have a markov process $P$ (on the joint space of the endogenous and exogenous states). We know from econometric theory that if we simulate a single time series using $P$ then asymptotically the distribution of this time series will convergence to what we are calling the stationary distribution $\mu = P \mu$.

Two minor points are needed to implement this. First, we have to start from somewhere, and we don't one the point we start from to matter, so we do 'burnin': simulate for $B$ periods, throw these away except for where we end up, and now use this as the starting point for simulating our distribution (because the starting point is random it will not matter on average; it is also likely to be a 'typical' point). Second, simulating a single time series is slow, so in practice we simulate lots of time series in parallel on the cpu; while this is not in line with the underlying theory justifying the approach it works just fine in practice.

The pseudocode for simulating the agent distribution follows. Note everything is done based on the index of $(a,z)$, not the value.
\begin{algorithmic}
 \State We will simulate $M$ observations. Done as $N$ simulations of length $T$ ($M=N \times T$)
 \For $i$ count from $1$ to $N$ (Parallel-for over cpu cores)
  \State Choose starting point $(a_0,z_0)$
  \State Create $Dist_i=zeros($\textit{size of $(a,z)$ space}$)$
  \For $t$ counts from $1$ to $B+1$ \Comment $B$ is number of periods for burnin
   \State Draw new $(a_t,z_t)$ from $P(a',z'|a_{t-1},z_{t-1})$
   \EndFor
  \For $t$ counts from $B+1$ to $T$ \Comment Same as during burnin, but now we will keep them
   \State Draw new $(a_t,z_t)$ from $P(a',z'|a_{t-1},z_{t-1})$
   \State $Dist_i(a_t, z_t)=Dist_i(a_t, z_t)+1$ (recall, using indexes, not values)
  \EndFor
  \State $Dist(:,i)=Dist_i$ \Comment Put the $Dist_i$ together
 \EndFor
 \State $Dist=sum(Dist,$in the $i$ dimension$)$
 \Comment $Dist$ is now a count of how many times each point in distribution was visited during the simulations
 \State $Dist=Dist/sum(Dist)$
 \Comment Normalize to be a probability distribution
 \\ \Return $Dist$
\end{algorithmic}
the implementation in VFI Toolkit uses the mid-point of the grids on $(a,z)$ as the initial starting point $(a_0,z_0)$ for all simulations (it will anyway be burned away). This can be controlled setting \textit{simoptions.seedpoint} (which has default value of the mid-point of the grids on $(a,z)$.

\subsubsection{Iterate Agent Distribution}
If we start from some initial guess for agent distribution $\mu_0$, and iterate $\mu_t=P \mu_{t-1}$, we know from the theory that this sequence of $\mu_t$ will converge to the stationary distribution $\mu$.\footnote{Typically $P$ is either a 'stochastic contraction' or, more likely for incomplete markets models, satisfies a 'monotone mixing condition' (or 'splitting condition', or 'monotone order-mixing condition') and either of these means the sequence will converge. Actually analytically proving either of these is non-trivial but has been done for some of the more basic/standard heterogenous agent incomplete market models.}

Iteration just repeatedly iterates until convergence, so psuedocode is just
\begin{algorithmic}
 \State Start from initial guess $\mu_{old}$
 \While {currdist$>$tolerance}
  \State $\mu=P \mu_{old}$ \Comment Iterate agent distribution
  \State $currdist=max(abs(\mu-\mu_{old}))$
  \State $\mu_{old}=\mu$ \Comment Update 'old'
 \EndWhile
 \\ \Return $\mu$ \Comment The distribution once convergence was reached
\end{algorithmic}
This just leaves the question of what to use as the initial guess $\mu_{old}$. By default VFI Toolkit just put equal weight on every grid point as the initial guess for $\mu_{old}$. You can give an initial guess for $\mu_{old}$ as \textit{simoptions.initialdist} (if guess is good, codes will run faster).\footnote{The toolkit used to use a small simulation as the initial guess, but the Tan improvement made the iteration so fast that getting a good initial guess was no longer worth the run time involved.}

The code implementing this contains one trick to speed it up, namely that it computes $\mu=P \mu_{old}$ multiple times before calculating $currdist$; because the distance calculation is non-trivial it is costly to compute it every step. This is controlled by \textit{simoptions.multiiter} which is set to 50 by default.

To avoid the risk of just iterating forever if the agent distribution is failing to converge, the codes impose a maximum number of iterations (it is an additional criterion of the while loop). This is controlled by \textit{simoptions.maxit} and is set to 50,000 by default.

You can control 'tolerance' by setting $simoptions.tolerance$.

\subsubsection{Tan Improvement to Iterating Agent Distribution}
By default VFI Toolkit will perform the 'Tan improvement' \cite*{Tan2020} to iteration which we now describe. In standard iteration we use the transition matrix $P$, which is the composition of the policy function $g(a'|a,z)$ and the transition matrix for the exogenous state $\pi_z(z'|z)$; $P=g \circ \pi_z$. The Tan improvement is the clever observation that actually instead of doing the iteration as $\mu=P \mu_{old}$, we can actually do it as a two-step process, with the first step being $\tilde{\mu}=\Gamma \mu_{old}$ and the second step is $\mu=\pi_z \tilde{\mu}$; where $\Gamma$ is just $g$ but as a transition from $(a,z)$ to $(a',z)$ rather than just to $a$ (notice that both $P$ and $\Gamma$ go from $(a,z)$, but $P$ goes to $(a',z')$ while $\Gamma$ goes to $(a',z)$). Conceptually this seems pointless, but computationally $\Gamma$ is very \textit{sparse} and so the Tan improvement is both vastly faster and uses less memory than standard iteration.

The pseudocode for iteration with the Tan improvement is thus,
\begin{algorithmic}
 \State Start from initial guess $\mu_{old}$
 \State Compute $\Gamma$ from $g$
 \While {currdist$>$tolerance}
  \State $\mu=\Gamma \mu_{old}$ \Comment First step of Tan improvement
  \State $\mu=\pi_z \mu$ \Comment Second step of Tan improvement
  \State $currdist=max(abs(\mu-\mu_{old}))$
  \State $\mu_{old}=\mu$ \Comment Update 'old'
 \EndWhile
 \\ \Return $\mu$ \Comment The distribution once convergence was reached
\end{algorithmic}
There is some reshaping of matrices involved in the Tan improvement that I have glossed over here, for a full explanation see \cite{Tan2020}. As in standard iteration our implementation does not check the $currdist$ every iteration as it is faster not to.

Because the Tan improvement makes the iteration step so much faster it is no longer worth generating a decent initial guess for $\mu_{old}$, and so the default just puts all the mass for the endogenous state on the mid-point of the grid, and for the exogenous state it iterates ten times with $\pi_z$ from putting equal mass on all grid points. You can give an initial guess for $\mu_{old}$ as \textit{simoptions.initialdist} (if guess is good, codes will run faster).



\subsection{Agent Distribution in Finite Horizon}
There are two ways to compute the agent distribution in finite horizon, simulation and iteration. The user must define/input the period-1 distribution. The iteration method is used by default as it has a better combination of speed and accuracy than simulation, but because the iteration requires more memory it will sometimes not be possible and simulation should then be used instead (default is \textit{vfoptions.iterate=1}, setting it to 0 means simulation will be used instead). When we use simulation we are essentially just doing a panel data simulation, but then converting the result into an agent distribution.

% Note that VFI Toolkit stores the finite agent distribution as a mass one probability distribution conditional on each age, as well as a vector of weights for ages.\footnote{Two motives. That every age-conditional distribution should be mass one is something we can check when trying to debug. I suspect this is more robust to numerical rounding error.}

\subsubsection{Simulating the Agent Distribution}
The pseudocode for simulating the agent distribution follows. Note everything is done based on the index of $(a,z)$, not the value. We use $\mu_j$ to denote the distribution at age/period $j$, and $g$ denotes policy function, $\pi_z$ denotes the transition matrix for the exogenous state.
\begin{algorithmic}
 \State Define initial age-1 distribution $\mu_1(a,z)$
 \State Define age-weights $\omega(j)$
 \State We will perform $N$ simulations of length $J$
 \For $i$ count from $1$ to $N$ (Parallel-for over cpu cores)
  \State Draw starting point $(a_1,z_1)$ from $\mu_1(a,z)$
  \State Create $Dist_i=zeros($\textit{size of $(a,z,j)$ space}$)$
  \For $j$ counts from $1$ to $J$
   \State Get $a_j=g(a_{j-1},z_{j-1},j-1)$
   \State Draw $z_j$ from $\pi_z(z_j|z_{j-1};j)$
   \State $Dist_i(a_j, z_j,j)=Dist_i(a_j, z_j,j)+1$ (recall, using indexes, not values)
  \EndFor
 \EndFor
 \State $Dist(:,:,:,i)=Dist_i$ \Comment Put the $Dist_i$ together
 \State $Dist=sum(Dist,$in the $i$ dimension$)$
 \Comment $Dist$ is now a count of how many times each point in distribution was visited during the simulations
 \State $Dist(:,:,j)=Dist(:,:,j)/sum(Dist(:,:,j))$, for each $j$
 \Comment Normalize conditional on age to be a probability distribution
 \State Multiply $\omega$ along the $j$ dimension of $Dist$.
 \\ \Return $Dist$
\end{algorithmic}

\subsubsection{Iterating the Agent Distribution}
To iterate the agent distribution in finite horizon we need an initial distribution defined by the user, $\mu_1$, and then construct the transition matrix on $(a,z)$ for each period $j$ denoted $P_j$. $P_j$ is constructed by combining the policy function for period $j$, $g_j$, with the transition matrix on exogenous shocks for period $j$, $\pi_{z,j}$.

We just iterate $j$ times,
\begin{algorithmic}
 \State Define initial age-1 distribution $\mu_1(a,z)$ (must be mass 1)
 \State Define age-weights $\omega(j)$
 \For $j$ counts from $1$ to $J-1$
  \State Construct $P_j$ from $g_j$ and $\pi_{z,j}$
  \State $\mu_{j+1}=P_j \mu_j$ \Comment Iterate agent distribution
 \EndFor
 \State Store all the $\mu_j$ in $\mu$
  \State Multiply $\omega$ along the $j$ dimension of $\mu$.
 \\ \Return $\mu$
 \end{algorithmic}
In the actual codes the storing of $\mu_j$ in $\mu$ is done inside the for-loop, and we just have a $\mu_{new}$ and $\mu_{old}$ rather than all of the $\mu_j$.

VFI Toolkit by default will actually use the Tan improvement when iterating on the agent distribution and we turn to that now. We note that $P$ here takes us directly from $(a,z)$ to $(a',z')$.

\subsubsection{Tan Improvement Iterating the Agent Distribution}
The Tan improvement is the observation that instead of using $P$ we can use $g$ and $\pi_z$ one-by-one. This bit of genius means we do exactly the same, but in a way that is faster and less memory demanding (people spent 30+ years using $P$ before Tan saw that you could split it in two!).

We just iterate $j$ times,
\begin{algorithmic}
 \State Define initial age-1 distribution $\mu_1(a,z)$ (must be mass 1)
 \State Define age-weights $\omega(j)$
 \For $j$ counts from $1$ to $J-1$
  \State Construct $\Gamma_j$ from $g_j$
  \State $\mu_{temp}=\Gamma_j \mu_j$ \Comment First step of Tan improvement
  \State $\mu_{j+1}=\pi_{z,j} \mu_{temp}$ \Comment Second step of Tan improvement
 \EndFor
 \State Store all the $\mu_j$ in $\mu$
  \State Multiply $\omega$ along the $j$ dimension of $\mu$.
 \\ \Return $\mu$
 \end{algorithmic}
There is some reshaping of matrices involved in the Tan improvement that I have glossed over here, for a full explanation see \cite{Tan2020}.

The Tan improvement can be understood as breaking the move from $(a,z)$ to $(a',z')$ into two steps. The first step is from $(a,z)$ to $(a',z)$, and the second step is from $(a',z)$ to $(a',z')$.

\subsubsection{Tan Improvement with Semi-Exogenous shocks}
Semi-exogenous shocks are exogenous states for which the transition probabilities depend on a decision variable (see Life-Cycle Model 30 for an example). So $\pi_{semiz}(semiz'|semiz,d)$ gives the probability of next-period semi-exogenous state based on the current semi-exogenous state and on the decision variable.

Using the Tan improvement requires a minor modification. Note that we cannot treat $semiz$ like we treat $z$, in the second step of the Tan improvement, because semiz depends on $d$ (and hence on $a$ and $z$). Instead we have to include $semiz$ in the first step of the Tan improvement.

We just iterate $j$ times,
\begin{algorithmic}
 \State Define initial age-1 distribution $\mu_1(a,semiz,z)$ (must be mass 1)
 \State Define age-weights $\omega(j)$
 \For $j$ counts from $1$ to $J-1$
  \State Construct $\Gamma_j$ from $g_j$ and $\pi_{semiz,j}$
  \State $\mu_{temp}=\Gamma_j \mu_j$ \Comment First step of Tan improvement
  \State $\mu_{j+1}=\pi_{z,j} \mu_{temp}$ \Comment Second step of Tan improvement
 \EndFor
 \State Store all the $\mu_j$ in $\mu$
  \State Multiply $\omega$ along the $j$ dimension of $\mu$.
 \\ \Return $\mu$
 \end{algorithmic}
The difference from the standard Tan improvement is all in the first step, which now includes the semi-exogenous state transitions alongside the standard endogenous state transitions.
There is some reshaping of matrices involved in the Tan improvement that I have glossed over here, for a full explanation see \cite{Tan2020}. Note that now our distribution is on $(a,semiz,z)$, and $\Gamma_j$ is used in the first step to change to $(a',semiz',z)$, and then the second step goes to $(a',semiz',z')$.\footnote{Getting the Tan improvement to work here does require the order $(a,semiz,z)$, in fact the whole toolkit got rewritten just to be able to do this as it originally had $z$ then $semiz$.}

\subsubsection{Iterating the Agent Distribution with 'TwoProbs' (Linear Interpolation)}
Sometimes our policy is 'between' grid points, e.g., while using experienceasset or \href{https://github.com/robertdkirkby/GridInterpolationLayer}{grid interpolation layer}. Previously our policy (for next period endogenous state) was always on the grid, which we can think of as a policy with a single next-period grid point to which we go with probability one. When the policy is 'between' grid points we will think about going to two grid points, those below and above our actual policy, with probabilities that depend on how close our policy is to each grid point. We call the two grid points the 'lower' grid point and the 'upper' grid point. Obviously the two probabilities should sum to one (we go to somewhere next period with certainty), and so we can just think of the probability of the upper grid point as one minus the probability of the lower grid point. Thus we just need to decide how to set the probability of the lower grid point, and we will use linear interpolation. Summing up so far we convert or 'between' grid point policy into a two-valued policy for the lower and upper grid points, together with two probabilities for these lower and upper grid points. We can then perform standard iteration with this 'twoprobs' policy.

So for example, say the policy value is $a*$ (we will have one policy value for each point in the state space, so we will be doing this many times, but here just consider one). Our grid of values on the endogenous state is $a\_{}grid={a_1,a_2,\dots, a_{n\_{}a}}$. We first need to figure out what is the lower grid point, the grid point immediately below $a*$. We get this by solving $\argmax_{i=1,...,n\_{}a} a_i<=a*$; the lower grid point is the maximum point in $a\_{}grid$ that is below $a*$. Call $a_{i_l}$ the lower grid point. By definition the upper grid point is one above this, so $a_{i_l+1}$. The probability of the lower grid point is then given by $\pi_{l}=1-\frac{a*-a_{i_l}}{a_{i_l+1}-a_{i_l}}$; this is linear interpolation, you can see that if $a*$ is one-quarter of the distance from the lower to the upper grid points, then the probability of the lower grid point is three-quarters, if it is halfway the probability is half. The probability of the upper grid point is then $1-\pi_{l}$; because the two probabilities must sum to one.

In this way we convert the policy function $Policy(a,z,j)$ into a lower grid point policy $Policy_{l}(a,z,j)$ and a lower grid point probability $Policy_{\pi,l}(a,z,j)$. We don't need to store the upper grid point policy as it is just the lower grid point policy plus one, and we similarly don't need to store the upper grid point probability as it is just one minus the lower grid point probability.

First, let's discuss what $P_j$, the transition matrix from period $j$ state to period $j+1$ state would look like if there were no exogenous markov $z$ variables, just an endogenous state $a$ (so this paragraph is a simpler setting than the rest of this section). This will help to illustrate the differences in $P_j$ when constructed from a 'between' grid point policy, rather than an on-grid policy. To iterate the agent distribution in finite horizon we need an initial distribution defined by the user, $\mu_1(a)$, and then construct the transition matrix on $(a)$ for each period $j$ denoted $P_j$. When $Policy(a,j)$ was a single-value, we got that $P_j$ was a matrix from $(a)$ to $(a')$, and in each row there was a single 1 (in the index for next period $a'$ that corresponds with the policy, given the row/index for this period $a$), together with many zeros. Now that $Policy(a,j)$ is 'between' grid points, we instead construct $P_j$ from the lower-grid point policy $Policy_l(a,j)$ and the lower grid point probability $Policy_{\pi,l}(a,j)$ (and the upper grid point policy and probability). Thus we now get that $P_j$ is again a matrix from $(a)$ to $(a')$, but now in each row there are two non-zeros corresponding to the indexes for the lower and upper grid points, and the values of these non-zeros are the probabilities of the lower and upper grid point, respectively.

Now, let's put the markov $z$ back into the problem, the on-grid point $P_j$ would now be $n\_{}z$ non-zeros per row, with probabilities taken from $pi\_{}z$, while the between grid point $P_j$ will be $2*n\_{}z$ non-zeros per row, with probabilities taken from $pi\_{}z \times [\pi_{l},1-\pi_{l}]$. To iterate the agent distribution in finite horizon we need an initial distribution defined by the user, $\mu_1(a,z)$, and then construct the transition matrix on $(a,z)$ for each period $j$ denoted $P_j$. When $Policy(a,z,j)$ is a single-value, we get that $P_j$ was a matrix from $(a,z)$ to $(a',z')$, and in each row there are $n\_{}z$ non-zeros (in the index for next period $a'$ that corresponds with the policy, given the row/index for this period $a$, crossed with the indexes for all next period $z'$), together with the probabilities for next period $z'$ (from $pi\_{}z(z_c,:)$, where $z_c$ is this period $z$ index that corresponds to the row of $P_j$. Now that $Policy(a,z,j)$ is 'between' grid points, we instead construct $P_j$ from the lower-grid point policy $Policy_l(a,z,j)$ and the lower grid point probability $Policy_{\pi,l}(a,z,j)$ (and the upper grid point policy and probability). Thus we now get that $P_j$ is again a matrix from $(a,z)$ to $(a',z')$, but now in each row there are $n\_{}z$ non-zeros corresponding to the indexes for the lower and upper grid points crossed with all the next-period $z$ indexes, and the values of these non-zeros are the probabilities of the lower and upper grid point crossed with the probabilities of next period $z'$ from $pi\_{}z(z_c,:)$, respectively.

Now that we have created $P_j$ for $j=1,\dots, J$, the rest of the algorithm for the agent distribution in a finite horizon is unchanged. Of course because $P_j$ contains more non-zeros the rest of the algorithm will take longer to compute,\footnote{Only the non-zeros matter in practice as you will always want to implement this using sparse matrices.} but the actual steps are identical.

So to repeat what we saw in the case with a policy on grid, but now with a 'different' $P_j$ because our policy was between grid points, we now just iterate $j$ times,
\begin{algorithmic}
 \State Define initial age-1 distribution $\mu_1(a,z)$ (must be mass 1)
 \State Define age-weights $\omega(j)$
 \For $j$ counts from $1$ to $J-1$
  \State Construct $P_j$ from $g_j$ and $\pi_{z,j}$
  \State $\mu_{j+1}=P_j \mu_j$ \Comment Iterate agent distribution
 \EndFor
 \State Store all the $\mu_j$ in $\mu$
  \State Multiply $\omega$ along the $j$ dimension of $\mu$.
 \\ \Return $\mu$
 \end{algorithmic}
In the actual codes the storing of $\mu_j$ in $\mu$ is done inside the for-loop, and we just have a $\mu_{new}$ and $\mu_{old}$ rather than all of the $\mu_j$.

VFI Toolkit by default will actually use the Tan improvement when iterating on the agent distribution, but once we understand how to create $P_j$ for 'two probs' from our between grid policy, that is not really changing anything from when we saw the Tan improvement earlier so it is not repeated here.

Note, the use of 'two probs' is not dependent on the use of linear interpolation. In VFI Toolkit it is (currently) always linear interpolation that is used to create the 'two probs' policies, but in principle any other approach to creating the two probabilities could be used.

One last point. VFI Toolkit has to create the 'two probs' $P_j$ here for a few different settings. One is experienceasset, in which case the Policy is for an $a*$ value as described above. But another is for grid interpolation layer, in which case the way to construct $Policy_l$ and $Policy_{\pi,l}$ is slightly different (but after we have these is unchanged). Specifically, for grid interpolation layer the $Policy$ contains both the lower grid point for $a\_{}grid$ ---so we already directly have $Policy_l$--- and the index for the 'interpolation layer' which is $n2$ points evenly spaced between two consecutive grid points of $a\_{}grid$ (the index can be 1 to $n2+2$ as it includes the two end points of the original grid). Since these interpolation layer points are evenly spaced, we can trivially get the probability of the lower grid point as just $1-($interpolation-layer-index minus one$)/(n2+1)$. The rest is as above. 'residualasset' is dealt with the same as 'experienceasset' as it too has Policy as a between grid point $a*$ value.

\subsubsection{Iterating the Agent Distribution with 'uProbs'}
When using riskyasset or experienceassetu, we get not one between-grid-points value as the policy for next period endogenous state, but $n\_{}u$ between-grid-points values (each value of the i.i.d. shock $u$ results in a different value). We can just create $Policy_l$ and $Policy_{\pi,l}$ for each of these $n\_{}u$ values following what we did in the previous section for 'two probs'. We then create $P_j$, which is now $2*n\_{}u*n\_{}z$ non-zero values per row (two for the upper and lower grid points, for each of the $n\_{}u$ values, and then crossed with the $n\_{}z$ next period values for the markov $z'$). From here the iteration of the agent distribution is standard, and things like the Tan improvement, or using semi-exogenous or i.i.d. exogenous states can all be done as usual.




\section{Stationary General Equilibrium}
To solve for a stationary general equilibrium, we are looking for some parameters that are determined in general equilibrium, call them $\theta$ such that our general equilibrium equations equal zero, $GEeqn(\theta)=0$.

There are various algorithms for doing this (\textit{heteroagentoptions.fminalgo} determines which is used), but all of them essentially boil down to the following,
\begin{algorithmic}
 \State Guess some initial $\theta^0$
 \State Define $GeneralEqmConditions=\infty$
 \While $GeneralEqmConditions>tolerance$
 \State Solve the value fn and policy (given current $\theta^k$)
 \State Solve the agent distribution (given current $\theta^k$ and policy we just found)
 \State Solve the aggregate variables at each $t$ (given current $\theta^k$ and policy/agent dist we just found)
 \State Evaluate $GeneralEqmConditions=GEeqn(AggVars,\theta^k)$ (given current $\theta^k$ and agg vars)
 \State Update $\theta^{k+1}$ (in some way, likely based on $\theta^k$ and $GeneralEqmConditions$)
 \EndWhile
 \\ \Return $\theta$
\end{algorithmic}
The difference between the different \textit{fminalgo} is how we update the $\theta$. Note that in almost all of the \textit{fminalgo} the while loop is being done by an optimization algorithm (like Matlab's \textit{fminsearch()}).

Whether the problem is finite or infinite-horizon value functions, and whether or not there are permanent types of different kinds of shocks is unimportant, in the sense that the above steps are all the same steps, just that the details of their implementation is changed.

When there are multiple general equilibrium equations, $GeneralEqmConditions$ is (by default) defined as the sum-of-squares of the individual general equilibrium equations.


\section{Transition Path General Equilibrium}
We consider a transition path of length $T$ time periods. Solving for a transition path general equilibrium, we are looking for some path-of-parameters that are determined in general equilibrium, call them $\Theta$ such that our general equilibrium equations equal zero in every time period, $GEeqn(\Theta)=0$. Note that at some level this is the same thing we did for a stationary equilibrium, except that now $\Theta=\{\theta_1,....,\theta_T\}$ is a time-path of general-eqm parameters, rather than just parameters at a single point in time.

There are various algorithms for doing this (\textit{heteroagentoptions.fminalgo} determines which is used), but all of them essentially boil down to the following,
\begin{algorithmic}
 \State Guess some initial $\Theta^0$ (time-path on general-equilibrium parameters)
 \State Define $GeneralEqmConditions=\infty$
 \While $GeneralEqmConditions>tolerance$
 \State Solve backward, $t=T-1,\dots,2,1$ for the value fn and policy at each $t$ (given current $\theta^k_t$)
 \State Solve forward, $t=1,2,\dots,T$ for the agent distribution (given current $\theta^k_t$ and $Policy_t$ we just found)
 \State Solve the aggregate variables (given current $\theta_t^k$ and policy/agent dist we just found for $t$)
 \State Evaluate $GeneralEqmConditions=[GEeqn_1(AggVars_1,\theta_1^k),$ $GEeqn_1(AggVars_1,\theta_1^k),$$\dots,$ $GEeqn_T(AggVars_T,\theta_T^k)]$ (evaluate general eqm eqns at each time period)
 \State Update $\Theta^{k+1}$ (in some way, likely based on $\Theta^k$ and $GeneralEqmConditions$)
 \EndWhile
 \\ \Return $\theta$
\end{algorithmic}
The difference between the different \textit{fminalgo} is how we update the $\Theta$. The most common way to update is as a shooting algorithm.

In a stationary equilibrium, parameters other than those parameters which were being determined in general equilibrium (the general equilibrium parameters) were all constant (in time) and kept in Parameters structure. When we solve transition paths, some parameters might vary over time, and this is done with the ParamPath (a path on exogenous parameters), while any other parameters not in this path are assumed to be constant at their values in Parameters structure. It is not described in algorithm above, but obviously we have to make sure to use the correct parameters in each time period.

It is important to solve backward for policy and forward for agent distribution. For convenience, aggregate variables and general eqm conditions are often solved forwards while doing agent distribution, but in principle how we deal with these (in terms of timing) is not important is most models (can be important if we use lagged/leading parameter values for something).

Whether the problem is finite or infinite-horizon value functions, and whether or not there are permanent types of different kinds of shocks is unimportant, in the sense that the above steps are all the same steps, just that the details of their implementation is changed.



\nocite{VFIToolkit}
\bibliographystyle{plainnat}
\bibliography{VFI_PseudoCodes_refs}

\end{document}
